\chapter{Discussion}
\label{chap:discussion}

\section{Summary of Findings}

The thesis pursued a single operational question: can a heterogeneous
knowledge graph built from Indian court judgments, filtered to be
structurally leakage-safe, support an outcome-prediction model whose
accuracy is attributable to legal reasoning rather than to identity or
decision shortcuts? The chapters above answer that affirmatively but
with important qualifications.

The pipeline produces a legally coherent corpus: the top-PageRank
entities recovered per bucket are exactly the statutes, provisions,
and courts a practitioner would expect, and the leakage-focused
pruning of Chapter~\ref{chap:kg} (dropping auxiliary NER types,
forbidding shortcut edges from identity proxies, keeping courts and
counsel local) is restrictive without being destructive: the resulting
graph is large, connected, and hub-dominated, with IPC, CrPC, and the
apex courts bridging all five legal domains.

The HGT learns cross-case structural signal that is not reducible to
its text features. The post-GNN probing classifier outperforms the
pre-GNN baseline at fixed capacity with a $16\times$ smaller
representation, and the t-SNE geometry shows that message passing
reorganises an already-strong text representation into a more
outcome-informative case embedding rather than replacing text with a
pure hub-traversal signal. The graph's contribution is best described
as structured aggregation of legal context, with cross-case shared
authorities providing a secondary, domain-dependent benefit.

The model resists identity shortcuts at four independent levels.
At the schema level, courts, judges, and lawyers are local and
shortcut edges from identity proxies are forbidden; at the text level,
rhetorical-role and phrase filtering remove direct decision leakage;
at the model level, the \texttt{no\_names} ablation matches the
baseline on every bucket; and at the representation level, probing,
saliency, and error analysis converge on the same conclusion. The
remaining difficult cases are the doctrinally crowded and document-heavy
judgments, not the de-identified ones.

Predictions are auditable case by case. The \path{Graph_Analyser}
pipeline of Chapter~\ref{chap:explainability} surfaces, for any test
case, the statutes, provisions, precedents, and argument spans the
frozen model relied on, and verbalises the bundle through a
Mistral-Small narrative. The same machinery exposes a structural
signature of confident error: the selected evidence can make the
predicted label look legally plausible while missing distinguishing
facts that explain the true outcome.

The cross-bucket representation also transfers, partially, to a fully
unseen legal domain. Evaluated zero-shot on
$11{,}769$ food-safety judgments
(Section~\ref{sec:cross_domain_food_safety}) the model is strictly
above the majority-class baseline with low across-fold variance,
which is evidence that the trained HGT carries a partially
domain-general legal-reasoning prior rather than memorising the
training domains' statutory hubs. The $\sim\!20$-point gap to
in-domain performance is substantial and bounds the strength of the
claim.

\section{Limitations}

\paragraph{Transductive evaluation.} The main Chapter~\ref{chap:results}
numbers are produced under a transductive protocol, in which the full
graph, including test cases' entity neighbourhoods, is visible
during training. This matches a realistic deployment setting (a new
judgment is appended to an existing canonical vocabulary), and the
food-safety experiment partially relaxes it by evaluating on an
entirely unseen legal domain. What the present evaluation does not
measure is finer-grained inductive generalisation within a familiar
domain, a calendar-year or court hold-out, where most entities are
known but the specific case neighbourhood is new. The per-year error
pattern suggests older cases are harder, which is a prior on how such
a split would behave.

\paragraph{LLM-generated outcome labels.} The win/loss labels are
produced by prompting Mistral on the OpenNyAI summary of each case
and are therefore noisy with respect to the underlying disposition.
The noise is not uniform: sexual-offences and financial-fraud matters
often have multi-part dispositions that do not map cleanly onto a
binary axis, and these are also the two worst-performing buckets,
at least partly a label-noise effect rather than a modelling
failure.

\paragraph{Leakage filtering is a lower bound, not a guarantee.} The
deterministic filters (rhetorical-role labels, field names, leakage
phrases) plus the structural argument (shared-entity types are law,
not identity) are neither complete: an RR classifier can mis-label a
decisive sentence, the phrase list is never exhaustive, and a
sufficiently expressive GNN could in principle reconstruct a decisive
fragment from non-decisive neighbours. The \texttt{no\_names} result
addresses one specific leakage mode, not all of them.

\paragraph{Small ablation margins and non-universal pooled optimum.}
The best configuration beats the baseline by $+1.7$ accuracy points
on the pooled dataset, with 5-fold error bars of similar magnitude,
so close ablations should be read as a cluster of comparable
configurations rather than a strict ranking. The per-dataset best
variant also differs by domain: family-matrimonial, motor-accidents,
and sexual-offences favour different ablations under accuracy or
macro-F1, so the pooled cross-bucket score is a useful compromise
over heterogeneous graph geometries rather than evidence that one
architectural setting is intrinsically best everywhere.

\section{Future Work}

Four directions follow most directly from the findings and
limitations above.

\paragraph{Stronger generalisation and leakage audits.} The natural
next experiments are inductive splits that the present transductive
protocol cannot evaluate: a calendar-year hold-out (e.g.\ train
pre-2020, test post-2020) and a court hold-out (train on all but one
High Court). The same hold-out machinery doubles as a sharper
leakage probe, a judge-hold-out and a court-hold-out would tighten
the identity-leakage claim from ``no evidence of shortcuts'' to
``shortcuts are quantitatively bounded''. The connectivity-score
diagnostic of Chapter~\ref{chap:graph_exploration} is the natural
metric for these splits.

\paragraph{Richer outcome targets.} The Mistral labeller already
produces a multi-label variant alongside the binary
\texttt{win}/\texttt{loss} target used here. Moving the classifier
head to a multi-label sigmoid or an ordinal regression over the
$[-1, +1]$ score would absorb the disposition ambiguity that
currently inflates label noise on sexual-offences and financial-fraud
matters and would enable disposition-specific error analysis
(\emph{remanded} versus \emph{dismissed}).

\paragraph{Cross-domain transfer and domain-adaptive heads.} The
food-safety experiment shows the cross-bucket HGT is a usable
zero-shot base on an unseen domain but also quantifies the
$\sim\!20$-point in-domain/out-of-domain gap. Two complementary
follow-ons map onto that finding: (i) light-weight fine-tuning of
the classifier head (with an optional \textsc{Case}-node adapter) on
a modest labelled target-domain sample, to measure how quickly the
gap closes, and (ii) a domain-conditioned head that reads a domain
embedding alongside the post-GNN case embedding so that one unified
model can be evaluated across training and held-out domains without
per-domain retraining. Scaling the bucket inventory to additional
domains (tax, labour, consumer) is a related engineering exercise
that the pipeline already supports.

\paragraph{Expert-in-the-loop validation of case-level explanations.}
The \path{Graph_Analyser} pipeline already emits, for any case, the
ranked statutes, provisions, precedents, and argument spans the frozen
HGT relied on. Sampling
predictions across confidence and difficulty strata and putting the
resulting evidence bundles in front of a legal expert would test
(i) whether the surfaced authorities are legally defensible reasons
for the predicted outcome, (ii) whether the selected argument evidence
captures the legally decisive facts, and (iii) whether confident-wrong
cases correspond to a coherent class of distinguishing-fact omissions.
That turns the explainer from an
internal diagnostic into an auditable artefact and is the cleanest
path from offline evaluation to a deployable review tool.

\medskip

Taken together, the contributions of this thesis are a
leakage-safe construction pipeline for legal knowledge graphs, an
HGT-based prediction model whose accuracy survives an aggressive
identity-leakage audit, a case-level interpretability pipeline that
exposes both correct reasoning and a structural signature of confident
error, and a zero-shot evaluation on a fully unseen legal domain that
quantifies the cross-domain transferability of the learned
representation.
